\section[Regra da Substituição (5.5)]{A Regra da Substituição (5.5) --- a regra da cadeia lida ao contrário}

\begin{definicao}[{Regra da Substituição}]
Se $u=g(x)$ é derivável e $f$ é contínua na imagem de $g$:
\[
\int f\big(g(x)\big)\,g'(x)\dx=\int f(u)\du,\qquad du=g'(x)\dx .
\]
Para integrais definidas, \textbf{troque também os limites}:
\[
\int_a^b f\big(g(x)\big)\,g'(x)\dx=\int_{g(a)}^{g(b)}f(u)\du .
\]
Por que funciona: se $F'=f$, então $\frac{d}{dx}F(g(x))=F'(g(x))g'(x)=f(g(x))g'(x)$ (regra da cadeia). A substituição só ``desfaz'' a cadeia. O $du=g'(x)\dx$ vem da notação de Leibniz: $\frac{du}{dx}=g'(x)$.
\end{definicao}

\begin{metodo}[{Como reconhecer e como executar}]
\textbf{Reconhecer:} há uma função ``de dentro'' $g(x)$ --- argumento de $\sen$, $\cos$, $e^{(\cdot)}$, $\ln$; base de uma potência; radicando; denominador --- \emph{e} a derivada $g'(x)$ aparece multiplicando, a menos de uma constante.

\textbf{Executar:}
\begin{enumerate}[leftmargin=1.6em, itemsep=1pt]
\item $u=g(x)$ (a função de dentro).
\item $du=g'(x)\dx$. Se falta só uma constante $k$, escreva $g'(x)\dx=du$ e ajuste: $\dx=\frac{du}{g'(x)}$.
\item Reescreva a integral \textbf{inteiramente} em $u$. \emph{Nenhum $x$ pode sobrar.} Se sobrar, a substituição está errada ou não é o método.
\item Integre em $u$ (tabela).
\item Indefinida: volte para $x$ e $+C$. Definida: use os limites novos $g(a)$, $g(b)$ e \emph{não} volte para $x$.
\item Verifique derivando (indefinida) ou pelo sinal/tamanho esperado (definida).
\end{enumerate}
\textbf{Não é substituição simples} quando o fator que sobra não é constante: $\int e^{x^2}\dx$ (falta um $x$), $\int x^2\cos(x^4)\dx$ (falta $x^3$, sobra $x^2$). Essas não têm primitiva elementar --- não invente uma.
\end{metodo}

\begin{exemplo}[{Exemplo 5.1 --- o exemplo-modelo (Stewart, 5.5, Exemplo 1)}]
$\displaystyle\int x^3\cos(x^4+2)\dx$.

\textbf{Classificação:} $\cos(\;\cdot\;)$ com argumento $x^4+2$, e a derivada $4x^3$ aparece a menos da constante $4$. Substituição.

$u=x^4+2$, \; $du=4x^3\dx$ \;$\Rightarrow$\; $x^3\dx=\dfrac{du}{4}$.
\[
\int x^3\cos(x^4+2)\dx=\int\cos u\cdot\frac{du}{4}=\frac14\int\cos u\du=\frac14\sen u+C=\frac14\sen(x^4+2)+C .
\]
\textbf{Verificação:} $\dfrac{d}{dx}\Big[\tfrac14\sen(x^4+2)\Big]=\tfrac14\cos(x^4+2)\cdot4x^3=x^3\cos(x^4+2)$ \ok.
\end{exemplo}

\begin{exemplo}[{Exemplo 5.2 --- raiz de função afim}]
$\displaystyle\int\sqrt{2x+1}\dx$. \quad $u=2x+1$, $du=2\dx$, $\dx=\dfrac{du}{2}$:
\[
\int u^{1/2}\frac{du}{2}=\frac12\cdot\frac{u^{3/2}}{3/2}+C=\frac13u^{3/2}+C=\frac13(2x+1)^{3/2}+C .
\]
Verificação: $\frac13\cdot\frac32(2x+1)^{1/2}\cdot2=(2x+1)^{1/2}$ \ok.
\end{exemplo}

\begin{exemplo}[{Exemplo 5.3 --- o $x$ do numerador é a pista}]
$\displaystyle\int\frac{x}{\sqrt{1-4x^2}}\dx$. \quad $u=1-4x^2$, $du=-8x\dx$ \;$\Rightarrow$\; $x\dx=-\dfrac{du}{8}$:
\[
\int u^{-1/2}\Big(-\frac{du}{8}\Big)=-\frac18\cdot\frac{u^{1/2}}{1/2}+C=-\frac14\sqrt u+C=-\frac14\sqrt{1-4x^2}+C .
\]
(Sem o $x$ no numerador seria $\arcsen(2x)/2$ --- outra integral, outra técnica. O fator $x$ decide.)
\end{exemplo}

\begin{exemplo}[{Exemplo 5.4 --- $\int\tg x\dx$ e o padrão $\int\frac{g'}{g}$}]
$\displaystyle\int\tg x\dx=\int\frac{\sen x}{\cos x}\dx$. \quad $u=\cos x$, $du=-\sen x\dx$:
\[
\int\frac{-du}{u}=-\ln|u|+C=-\ln|\cos x|+C=\ln|\sec x|+C .
\]
Padrão geral: $\displaystyle\int\frac{g'(x)}{g(x)}\dx=\ln|g(x)|+C$. Exemplos: $\int\frac{2x}{x^2+1}\dx=\ln(x^2+1)+C$; $\int\frac{e^x}{e^x+3}\dx=\ln(e^x+3)+C$.
\end{exemplo}

\begin{exemplo}[{Exemplo 5.5 --- definida, trocando os limites (Stewart, 5.5, Exemplo 7)}]
$\displaystyle\int_1^2\frac{dx}{(3-5x)^2}$. \quad $u=3-5x$, $du=-5\dx$, $\dx=-\dfrac{du}{5}$. Limites: $x=1\Rightarrow u=-2$; $x=2\Rightarrow u=-7$.
\begin{align*}
\int_1^2\frac{dx}{(3-5x)^2}&=-\frac15\int_{-2}^{-7}u^{-2}\du=-\frac15\Big[-\frac1u\Big]_{-2}^{-7}
=-\frac15\Big[\Big(-\frac{1}{-7}\Big)-\Big(-\frac{1}{-2}\Big)\Big]\\
&=-\frac15\Big[\frac17-\frac12\Big]=-\frac15\cdot\Big(-\frac{5}{14}\Big)=\frac{1}{14}.
\end{align*}
\textbf{Conferindo:} em $[1,2]$ o integrando é positivo e vai de $\frac14$ a $\frac1{49}$, então a integral está entre $\frac1{49}\approx0{,}02$ e $\frac14$: $\frac1{14}\approx0{,}071$ \ok. Note que os limites em $u$ ficaram ``ao contrário'' ($-2$ para $-7$) --- isso é normal; não os reordene.
\end{exemplo}

\begin{exemplo}[{Exemplo 5.6 --- $\ln$ como função de dentro}]
$\displaystyle\int_1^{e}\frac{\ln x}{x}\dx$. \quad $u=\ln x$, $du=\dfrac{dx}{x}$; $x=1\Rightarrow u=0$; $x=e\Rightarrow u=1$:
\[
\int_0^1 u\du=\Big[\frac{u^2}{2}\Big]_0^1=\frac12 .
\]
\end{exemplo}

\begin{definicao}[{Simetria (Teorema 7 da seção 5.5) --- economiza conta e evita erro}]
Se $f$ é contínua em $[-a,a]$:
\[
f\ \text{par}\ (f(-x)=f(x))\ \Rightarrow\ \int_{-a}^{a}f\dx=2\int_0^a f\dx;
\qquad
f\ \text{ímpar}\ (f(-x)=-f(x))\ \Rightarrow\ \int_{-a}^{a}f\dx=0 .
\]
Geometria: função par tem gráfico simétrico ao eixo $y$ (duas áreas iguais); ímpar tem simetria pela origem (as áreas se cancelam). Regras rápidas: par $\times$ par $=$ par; ímpar $\times$ ímpar $=$ par; par $\times$ ímpar $=$ ímpar; $x^{\text{par}}$, $\cos$, $|x|$ são pares; $x^{\text{ímpar}}$, $\sen$, $\tg$ são ímpares.
\end{definicao}

\begin{exemplo}[{Exemplo 5.7 --- simetria (Stewart, 5.5, Exemplos 9 e 10)}]
\textbf{(a)} $\displaystyle\int_{-2}^{2}(x^6+1)\dx=2\int_0^2(x^6+1)\dx=2\Big[\frac{x^7}{7}+x\Big]_0^2=2\Big(\frac{128}{7}+2\Big)=\frac{284}{7}$.

\textbf{(b)} $\displaystyle\int_{-1}^{1}\frac{\tg x}{1+x^2+x^4}\dx=0$: numerador ímpar, denominador par, quociente ímpar. Sem calcular nada.
\end{exemplo}

\begin{armadilha}[{Onde a substituição dá errado}]
\begin{itemize}[leftmargin=1.4em, itemsep=1pt]
\item Sobrar $x$ na integral em $u$ (ex.: escrever $\int x\cos u\du$). Ou você resolve o $x$ em função de $u$, ou o método não é este.
\item Esquecer o fator de ajuste: $\int\cos(5x)\dx=\frac15\sen(5x)+C$, não $\sen(5x)+C$. (Derive e veja o $5$ aparecer.)
\item Na definida, misturar: voltar para $x$ \emph{e} usar os limites em $u$, ou ficar em $u$ e usar os limites em $x$. Escolha um caminho.
\item Escolher $u$ = ``a expressão toda'' em vez da função de dentro. Em $\int x^3\cos(x^4+2)\dx$, $u$ é $x^4+2$, não $\cos(x^4+2)$.
\item Esquecer o módulo em $\ln|u|$ quando $u$ pode ser negativo.
\end{itemize}
\end{armadilha}

\begin{exercicios}[{Exercícios similares --- Parte 5}]
\begin{enumerate}[leftmargin=1.6em, itemsep=2pt, label=\textbf{5.\arabic*}]
\item $\displaystyle\int x^2e^{x^3}\dx$
\item $\displaystyle\int\frac{(\ln x)^2}{x}\dx$
\item $\displaystyle\int x\sqrt{x^2+1}\dx$
\item $\displaystyle\int_0^1 xe^{x^2}\dx$
\item $\displaystyle\int_0^{\pi/2}\cos x\,\sen(\sen x)\dx$
\item $\displaystyle\int_0^{\sqrt\pi}x\cos(x^2)\dx$ --- interprete o resultado.
\item $\displaystyle\int_{-\pi/4}^{\pi/4}\big(x^3+x^4\tg x\big)\dx$
\item $\displaystyle\int\frac{\sec^2x}{1+\tg x}\dx$
\end{enumerate}
\end{exercicios}
